\section{One generation from \texorpdfstring{$\mathbb C\otimes\HH\otimes\OO$}{C x H x O}}\label{sec:generation}

The algebra \(\mathbb C\otimes\HH\otimes\OO\), often called the Dixon algebra, has a long history as a candidate carrier of the
Standard Model: octonionic quark structure~\cite{gunaydin}, Dixon's programme~\cite{dixon}, charge quantization from a number
operator on \(\mathbb C\otimes\OO\)~\cite{furey2015}, the Standard Model gauge group as a symmetry of ladder operators on
\(\mathbb C\otimes\HH\otimes\OO\)~\cite{furey2018}, braided realizations~\cite{gresnigt}, and triality as the origin of three
generations~\cite{boyle}.  This section shows that the building blocks of Sections~\ref{sec:kernel} and~\ref{sec:octonions} reproduce
the content of one generation.  Most of the pattern is known; what is specific to this framework is listed in
\S\ref{subsec:gen-new}.  All statements were checked by explicit matrices on \(\mathbb C^{32}\), the anomaly traces in exact fractions
(Companion layers v4.9.0-r34, r35, r41--r43, pre-registered tests H19, H20, H26--H28).

\subsection{Conventions}

\(L_a\) and \(R_a\) denote left and right multiplication by \(e_a\) on \(\OO\) (Cayley--Dickson convention of the Companion), with
\(e_7e_4=-e_3\), \(e_7e_5=-e_2\), \(e_7e_6=e_1\).  Colour is \(SU(3)_c=\operatorname{Stab}_{G_2}(e_7)\).  Two complex units occur and must be
kept apart: the \emph{outer} unit \(i\), the \(\mathbb C\) factor, and the \emph{inner} unit \(L_{e_7}\) on \(\OO\).  Time reversal \(T\) is complex
conjugation of the \(\mathbb C\) factor.

\subsection{Charges from a number operator}

Following~\cite{furey2015}, put
\[
\alpha_a=\tfrac12\bigl(L_{e_a}+i\,L_{e_7e_a}\bigr)\quad(a=4,5,6),\qquad N=\sum_a\alpha_a^\dagger\alpha_a,\qquad Q=N/3 .
\]

\begin{proposition}[Charge sectors]\label{prop:charges}
The \(\alpha_a\) satisfy \(\alpha_a^2=0\) and \(\{\alpha_a,\alpha_b^\dagger\}=\delta_{ab}\), so \(N\) has eigenvalues \(0,1,2,3\) with
multiplicities \(1,3,3,1\) on \(\mathbb C\otimes\OO\cong\mathbb C^8\).  Colour commutes with \(Q\).  The sectors \(Q=0\) and \(Q=1\) are colour
singlets spanned by \(1\pm i\,e_7\); the sectors \(Q=\tfrac13\) and \(Q=\tfrac23\) are mutually conjugate colour triplets.  Moreover
\[
Q=\tfrac12-\tfrac{i}{6}\,L_{e_7}-\tfrac{i}{3}\,R_{e_7},\qquad TQT^{-1}=1-Q .
\]
\end{proposition}

\begin{proof}
The anticommutators follow from the Clifford relations \(L_aL_b+L_bL_a=-2\delta_{ab}\) on \(\OO\); a fermionic number operator on three
modes has the stated spectrum.  Colour invariance, the sector bases and the formula for \(Q\) were verified by explicit matrices.
\(T\) sends \(\alpha_a\) to \(-\alpha_a^\dagger\), hence \(N\) to \(3-N\).
\end{proof}

The pattern \(0,\tfrac13,\tfrac23,1\) with singlet, triplet, conjugate triplet, singlet is Furey's; in her reading the sectors are
\(\nu,\bar d,u,e^+\).  The multiplicities are forced by the algebra and are not a prediction.  Two remarks belong to this framework.
First, charge needs both units and both left and right multiplication; the scalar \(it\) of the kernel alone generates only a global
phase.  Second, the inner frame and the outer frame are related by \(P_\pm=\tfrac12(1\mp iL_{e_7})\), and an earlier statement of the
Companion that \(\operatorname{span}(e_1,e_2,e_3)\) carries the conjugate colour triplet confused the two (layer v4.9.0-r35).

\subsection{Weak doublets}

Let the Lorentz group act on \(\mathbb C\otimes\HH\) from the left (rotations \(L_q\), boosts \(iL_q\)), the weak \(SU(2)\) from the right
(\(T_3=-\tfrac i2R_k\)), colour on \(\OO\), and put \(Y/2=Q-\tfrac12\), \(Q_{\rm em}=T_3+Y/2\).

\begin{proposition}[Doublets]\label{prop:doublets}
The weak \(SU(2)\) commutes with the Lorentz group, colour and \(Q\).  The joint eigenspaces of \(Q\) and \(T_3\) carry
\[
\begin{array}{llll}
\toprule
Q\text{-sector} & Y/2 & Q_{\rm em}\ (T_3=\pm\tfrac12) & \text{reading}\\
\midrule
0\ (\text{singlet}) & -\tfrac12 & 0,\,-1 & (\nu,e)\\
\tfrac23\ (\text{triplet}) & +\tfrac16 & \tfrac23,\,-\tfrac13 & (u,d)\\
\tfrac13\ (\text{conj.\ triplet}) & -\tfrac16 & \tfrac13,\,-\tfrac23 & (\bar d,\bar u)\\
1\ (\text{singlet}) & +\tfrac12 & 1,\,0 & (e^+,\bar\nu)\\
\bottomrule
\end{array}
\]
and \(TQ_{\rm em}T^{-1}=-Q_{\rm em}\).  There is no weak singlet: \(T_1^2+T_2^2+T_3^2=\tfrac34\) everywhere.
\end{proposition}

The offset \(-\tfrac12\) is chosen here; it is fixed below.  Without a further structure every state is a weak doublet, so half of a
generation is missing.

\subsection{Chirality from a compensator}

The missing half is supplied by a projector of the form of the Heisenberg compensator (Definition~\ref{def:compensator}), now with
\(J\) the product of the two units:
\[
C_{\rm ch}=\tfrac12(1+J),\qquad J=i\,L_{e_7},\qquad T_a'=C_{\rm ch}T_a,\qquad
\tfrac Y2=Q-\tfrac12+(1-C_{\rm ch})T_3 .
\]

\begin{proposition}[One generation]\label{prop:generation}
\(J^2=1\), and \(C_{\rm ch}\) is a Hermitian projector of rank \(16\) commuting with colour, the Lorentz group, the right action of \(\HH\) and
\(Q\); its image is the sum of the sectors \(Q=0\) and \(Q=\tfrac23\).  Under \(SU(3)_c\times SU(2)'\times U(1)_Y\) the space decomposes into
\[
(3,2)_{1/6}\oplus(1,2)_{-1/2}\oplus(\bar3,1)_{-2/3}\oplus(\bar3,1)_{1/3}\oplus(1,1)_{1}\oplus(1,1)_{0},
\]
that is \((u,d)_L\), \((\nu,e)_L\), \(u^c\), \(d^c\), \(e^c\), \(\nu^c\): sixteen Weyl states with two Lorentz components each, and
\(Q_{\rm em}=T_3'+Y/2\).  The traces \(\operatorname{Tr}Y\), \(\operatorname{Tr}Y^3\), \(SU(2)^2U(1)\) and \(SU(3)^2U(1)\) vanish.  Finally
\(TC_{\rm ch}T^{-1}=1-C_{\rm ch}\).
\end{proposition}

The weak \(SU(2)'\) acts only on the image of \(C_{\rm ch}\); in the model this is chirality.  The frame change \(T\) exchanges image
and kernel, so matter frame and tachyon frame are, in the model, particle and antiparticle with mirrored chirality.  This is a reading.

\subsection{The choice is fixed by anomaly cancellation}

\begin{proposition}[Uniqueness up to the frame]\label{prop:anomaly}
Consider chirality projectors \(C=\tfrac12(1+J)\) on \(\mathbb C\otimes\OO\) with \(J\) Hermitian, \(J^2=1\), commuting with colour and \(Q\), of
rank \(16\), whose image is one colour singlet and one colour triplet; set \(T_a'=CT_a\), \(Q_{\rm em}=T_3+Q+\sigma\) and
\(Y/2=Q+\sigma+(1-C)T_3\).  Then:
\begin{enumerate}[label=(\alph*)]
\item There are exactly four such projectors, \(J\in\{\pm iL_{e_7},\pm iR_{e_7}\}\).
\item For each, \(\operatorname{Tr}Y=0\) holds exactly for \(\sigma=-\tfrac12\).
\item At \(\sigma=-\tfrac12\) all four anomaly traces vanish exactly for \(J=\pm iL_{e_7}\); for \(J=\pm iR_{e_7}\) one finds
\(\operatorname{Tr}Y^3=\pm\tfrac32\).
\item \(T(iL_{e_7})T^{-1}=-iL_{e_7}\).
\end{enumerate}
\end{proposition}

\begin{proof}
Each \(Q\)-sector is colour-irreducible with its own charge, so the commutant of colour and \(Q\) consists of scalars on the four
sectors; a Hermitian involution there is a choice of sign per sector, and the rank and image conditions leave four.  The traces
were computed in exact fractions.  (d) holds because \(T\) conjugates \(i\) and fixes \(L_{e_7}\).
\end{proof}

Under the assumptions of the setting nothing is left to choose: the gravitational anomaly fixes the offset, the remaining anomalies
select left against right multiplication, and what remains is the Standard Model content and its frame mirror.  The setting itself
--- weak \(SU(2)\) from the right, the form of \(Q_{\rm em}\), a projector on \(\mathbb C\otimes\OO\) --- stays an assumption.  That the
Standard Model is anomaly-free is of course known; the statement is only that, in this structure, anomaly cancellation alone decides
the choice.

\subsection{What is new, and what is not}\label{subsec:gen-new}

\emph{Known.}  The charge pattern from a number operator~\cite{furey2015}; the Standard Model gauge group and one generation in
\(\mathbb C\otimes\HH\otimes\OO\)~\cite{dixon,furey2018,gresnigt}; triality as a carrier of three generations~\cite{boyle}.

\emph{Specific to this framework.}  (i) The roles are placed on the two sides of \(\HH\) (Lorentz left, weak right) and on the frames.
(ii) Chirality comes from a projector of the same form \(\tfrac12(1+J)\) as the kernel's compensator, with \(J\) the product of the outer and
the inner unit.  (iii) Among the four admissible chirality projectors, anomaly cancellation alone selects the left one, up to the frame
change.  (iv) The frame change acts as charge conjugation with mirrored chirality.  We have not found (ii) and (iii) in this form in the
literature, but a systematic comparison is still to be done.

\emph{Not addressed.}  Why the weak \(SU(2)\) acts from the right; symmetry breaking; masses and mixing angles; why three generations and
not four (Section~\ref{sec:octonions}).  The construction was made under choices that were partly seen in advance, and this is disclosed in
the pre-registration notes.

\medskip\noindent\emph{Status.}  Propositions~\ref{prop:charges}--\ref{prop:anomaly} are established by proof or by explicit exact
computation under the stated setting.  The identification of the states with physical particles is the standard one; its derivation
from the kernel is not claimed.
